\section{Background and Related Work}\label{sec:related}

This section first describes the data, the evaluation roles, and the metrics
used in this study, and then reviews the related work.

\subsection{Detectors and data}\label{sec:data}
Table~\ref{tab:data} lists the evaluation cells, that is, each combination of
data, detector, and tracker (the trackers are introduced in Sect.~\ref{sec:setup}). On MOT17 val-half
\cite{milan2016mot16} the trackers use the published YOLOX-X detections
\cite{ge2021yolox}. The emission floor, the lowest detector score passed on,
is 0.01 or 0.1 (Table~\ref{tab:mot17}). On KITTI tracking training \cite{geiger2012kitti}, the trackers run
with two detectors that differ from the MOT17 detector: YOLOv8n
\cite{jocher2023yolov8} and RT-DETR-L \cite{zhao2024rtdetr}. To isolate the
calibration problem, the published MOT17 detections were also transformed
before the tracker by four monotone maps: cubing ($s^3$), halving ($0.5\,s$),
and logit temperatures $T=2$ and $T=0.5$, $\sigma(z/T)$. The trackers kept
their published thresholds. The external systems use the authors' released
detections on MOT17 val-half.

\begin{table*}[t]
\caption{Evaluation cells and their role. Development cells were visible while
the layer was designed; predeclared cells were run once with the frozen
policy.}
\label{tab:data}
\centering\footnotesize
\setlength{\tabcolsep}{4pt}
\begin{tabular}{@{}>{\raggedright\arraybackslash}p{36mm}>{\raggedright\arraybackslash}p{32mm}>{\raggedright\arraybackslash}p{50mm}>{\raggedright\arraybackslash}p{30mm}@{}}
\toprule
Data & Detector & Trackers & Role\\
\midrule
MOT17 val-half (7 sequences) & published YOLOX-X detections & ByteTrack (2 settings), OC-SORT, BoostTrack, SparseTrack & development\\
KITTI tracking training (21 sequences) & RT-DETR-L; YOLOv8n & ByteTrack and BoT-SORT (RT-DETR-L); OC-SORT (YOLOv8n) & development (detector transfer)\\
MOT17 val-half, score transforms & published YOLOX-X, transformed scores & ByteTrack (2 settings), OC-SORT, BoostTrack & development (robustness)\\
MOT17 val-half & released PD-SORT / Hybrid-SORT detections & PD-SORT, Hybrid-SORT & predeclared external\\
\bottomrule
\end{tabular}
\end{table*}


\subsection{Evaluation roles}\label{sec:roles}
Every cell has one of two roles (Table~\ref{tab:data}). Development cells
were visible while \scl{} was designed: MOT17 with four trackers, KITTI
with two other detectors, and the score transforms. The predeclared external
systems, PD-SORT and Hybrid-SORT, were selected and declared before they were
run, and they were run once with the frozen configuration. For each external
host, the reproduction against the authors' number is classed as EXACT
(within 0.2 HOTA) or CLOSE (within 1.0 HOTA with a documented, untuned
cause).

\subsection{Metrics and statistics}\label{sec:metrics}
Results are reported as HOTA \cite{luiten2021hota}, MOTA
\cite{bernardin2008clear}, and IDF1 \cite{ristani2016idf1}. They are computed
with the TrackEval library (commit 12c8791) and the official ground truth of
each benchmark, following MOTChallenge practice
\cite{dendorfer2021motchallenge}. On KITTI tracking, HOTA is the official
value averaged over cars and pedestrians.

All comparisons are paired. Differences are pooled over sequences, and their
95\% intervals come from a percentile bootstrap over sequences
\cite{efron1993bootstrap} with 10,000 resamples and seed 42. A difference is
called significant when its interval excludes zero. As a distribution-free
complement, per-sequence wins, ties ($|\Delta\mathrm{HOTA}|<0.01$), and losses
are also counted.

\subsection{Related work}\label{sec:lit}

Four research directions are related to \scl{}. Tracking-by-detection methods provide the
trackers that \scl{} serves. Work on detector confidence explains why their
fixed thresholds break when the detector changes. Adaptive thresholds inside
trackers and adaptive inference are the closest alternatives to \scl{}.

\subsubsection{Tracking by detection}
SORT predicts each box with a Kalman filter and matches the predictions to new
detections by overlap with the Hungarian method \cite{bewley2016sort}.
ByteTrack introduced a second association stage that uses low-confidence
detections to continue existing tracks \cite{zhang2022bytetrack}, and many
later trackers kept this two-stage design. Others, such as OC-SORT
\cite{cao2023ocsort}, use a single confidence threshold. All of these
trackers read detector scores through fixed thresholds that were tuned for
one detector, usually YOLOX \cite{ge2021yolox}. Seven published trackers of
both kinds are used in this paper as hosts and are kept unchanged
(Sect.~\ref{sec:setup}).

\subsubsection{Detector confidence and calibration}
Detectors of different families, such as YOLOX \cite{ge2021yolox}, YOLOv8
\cite{jocher2023yolov8}, YOLOv9 \cite{wang2024yolov9}, DETR
\cite{carion2020detr}, and RT-DETR
\cite{zhao2024rtdetr}, produce different score distributions. Their confidence
is not calibrated across models and conditions, and calibration methods for
object detection fit a mapping on labelled data \cite{kuppers2020detcalib}.
A tracker tuned for one detector therefore cannot assume that another
detector's scores mean the same thing, and labels to recalibrate them are
rarely available after deployment. Otsu's method selects a threshold from a histogram without labels by
maximizing the separation between two classes \cite{otsu1979}. \scl{}
applies it to the logits (log-odds) of the detector scores of recent frames.

\subsubsection{Adaptive thresholds inside trackers}
Several trackers adapt their own confidence thresholds. Adaptive ByteTrack
derives the confidence threshold of ByteTrack from the detections of each
frame \cite{ma2023adaptivebytetrack}. AdapTrack adapts the thresholds of its
matching step \cite{shim2024adaptrack}. ACT sets the high and low thresholds
of each frame from the first quartile of the confidence scores and the number
of candidates \cite{act2026adaptive}. These methods are built into one tracker
and use statistics of the current frame. \scl{} differs in
four ways. Its statistics come from past frames. It decides for the whole
stream whether to intervene at all, and otherwise keeps the tracker's own
operating point. It acts through a declared host contract on several
published trackers without modifying them. It is evaluated as one frozen
policy under confidence shift and on external systems.

\subsubsection{Adaptive inference}
Dynamic neural networks adapt their computation to each input
\cite{han2022dynamic}, and Chameleon re-profiles the configuration of video
analytics pipelines over time \cite{jiang2018chameleon}. These methods learn or
profile a policy for each model. Our earlier work, AC-MOT \cite{ismail2026acmot},
is a training-free controller that sets the detector operating point (score
threshold, NMS threshold, and input resolution) online from image and track
cues. Such a controller is tied to the detector it was calibrated with,
because a change of the detector's score scale leaves the image unchanged.
\scl{} reads the score stream instead and acts on the
tracker's input and thresholds. Table~\ref{tab:related} summarizes the
differences.

In summary, existing methods either adapt thresholds inside one tracker from
the current frame, calibrate scores with labelled data, or learn a policy for
one model. A training-free layer that calibrates itself from the detector's
own score stream, serves many unmodified trackers, and keeps a well-matched
tracker at its own accuracy is missing. \scl{} is designed to fill this gap.

\begin{table}[t]
\caption{Adaptive methods closest to \scl{} (BT: ByteTrack).}
\label{tab:related}
\centering\footnotesize
\setlength{\tabcolsep}{3pt}
\begin{tabular}{@{}>{\raggedright\arraybackslash}p{17mm}>{\raggedright\arraybackslash}p{18mm}>{\raggedright\arraybackslash}p{15mm}>{\raggedright\arraybackslash}p{18mm}@{}}
\toprule
Method & Adapted & Statistics & Where\\
\midrule
Adapt.\ BT \cite{ma2023adaptivebytetrack} & confidence threshold & current frame & inside BT\\
AdapTrack \cite{shim2024adaptrack} & matching thresholds & current frame & inside the tracker\\
ACT \cite{act2026adaptive} & high and low thresholds & current frame & inside the tracker\\
K{\"u}ppers \cite{kuppers2020detcalib} & detector scores & labelled data & detector output\\
Chameleon \cite{jiang2018chameleon} & pipeline configuration & profiling & analytics pipeline\\
\scl{} (this paper) & candidates, scores, thresholds & past frames, no labels & between detector and tracker\\
\bottomrule
\end{tabular}
\end{table}
